\subsection{An exact factor sieve in the first two original shells}

Keep \(p=12h+1\) prime, \(h\ge1\), and put
\(a_3=3h+1\), \(a_7=3h+2\). These are the original shells
with residuals \(R_{a_3}=3\) and \(R_{a_7}=7\).
They both lie in \(A_p\), and both are at most \(6h=(p-1)/2\).
For either retained integer \(a\), write its actual prime factorization
\(a=\prod_{\ell\mid a}\ell^{v_\ell}\). The exponent domain throughout
is exactly
\[
 \mathcal B_a=\prod_{\ell\mid a}\{0,\ldots,2v_\ell\},\qquad
 u(e)=\prod_{\ell\mid a}\ell^{e_\ell}.
\]
Unique prime factorization proves this is a bijection onto \(\Div(a^2)\),
with inverse \(e_\ell=v_\ell(u)\). Every prime, exponent cap and
integer \(a\) remains attached to the coordinates below.

\begin{theorem}[The complete residual-three factor sieve]
\label{thm:r3-factor-sieve}
At \(a=a_3\), define
\[
 P_1=\prod_{\substack{\ell\mid a\\\ell\equiv1\ (3)}}(2v_\ell+1),
 \qquad
 P_2=\prod_{\substack{\ell\mid a\\\ell\equiv2\ (3)}}(2v_\ell+1).
\]
Empty products are one. The full exterior and middle sets coincide,
and their complete exponent fibres are
\[
 E_a=M_a=
 \left\{u(e):e\in\mathcal B_a,\quad
           \sum_{\ell\equiv2\ (3)}e_\ell\equiv1\pmod2\right\}.
\]
In particular
\[
 |E_a|=|M_a|=\frac{P_1(P_2-1)}2,
 \qquad |E_a|+\frac{|M_a|}{2}=\frac{3P_1(P_2-1)}4.
\]
The shell has a witness if and only if \(a\) has a prime divisor
\(\ell\equiv2\pmod3\). Every such factor gives both channels
with primitive layer \(A=1\), via
\[
 u=a\ell,\quad A=1,\quad B=\ell,\quad D=a/\ell,
 \qquad C_0=(1+p\ell)/3,\quad C_1=(1+\ell)/3.
\]
Here the original tag \(\varepsilon=0\) uses \(C_0\), and
\(\varepsilon=1\) uses \(C_1\); their respective ordered triples are
\[
 (a,C_0D,p\ell C_0D),\qquad
 (a,pC_1D,p\ell C_1D).
\]
\end{theorem}
\begin{proof}
The identity \(a=3h+1\) gives \(a\equiv1\pmod3\), so no prime
factor is three. The exterior condition \(3\mid4u+1\) and the
middle condition \(3\mid u+a\) both say \(u\equiv2\pmod3\).
Every factor \(\ell\equiv1\pmod3\) contributes one to the
residue product, and every \(\ell\equiv2\pmod3\) contributes
\((-1)^{e_\ell}\). This proves both inclusions in the displayed
fibre equality, with the stated valuation inverse.

There are \(P_2\) exponent choices on the second residue class.
The number of even sums minus the number of odd sums equals
\[
 \prod_{\substack{\ell\mid a\\\ell\equiv2\ (3)}}
       \left(\sum_{e=0}^{2v_\ell}(-1)^e\right)=1,
\]
because every inner alternating sum is one. Thus the odd sums
number \((P_2-1)/2\); the first residue class independently
contributes all \(P_1\) choices. This proves the cardinality.
The complement \(u\mapsto a^2/u\) preserves the middle set:
its residue is \(1/2=2\) modulo three. Its only positive fixed
point would be \(u=a\), which has residue one and is absent.
Hence \(|M_a|\) is even, and the displayed original weighted
shell count is the stated integer.

If there is no second-class factor, every exponent sum is zero
and the fibre is empty. If \(\ell\equiv2\pmod3\) divides \(a\),
the integer \(u=a\ell\) has exponent \(v_\ell+1\le2v_\ell\)
at that prime and exponent \(v_t\) at every other prime \(t\).
It belongs to \(\Div(a^2)\) and has residue two. Its coordinates
satisfy \(a=ABD\), \(u=B^2D\), and \(\gcd(A,B)=1\) exactly.
Both \(C_0,C_1\) are positive integers since \(p\equiv1\pmod3\)
and \(\ell\equiv2\pmod3\). For each tag one has
\(3C_\varepsilon=1+p^{1-\varepsilon}\ell\) and
\(4a=p+3\). Multiplying gives
\[
 4AB C_\varepsilon D
   =A+p^{1-\varepsilon}B+pC_\varepsilon.
\]
Putting the three indicated reciprocals over \(pAB C_\varepsilon D\)
therefore proves their sum is \(4/p\). This proves the complete
arithmetic return with the original ordering and first denominator.
\end{proof}

\begin{theorem}[The complete residual-seven factor sieve]
\label{thm:r7-factor-sieve}
At \(a=a_7\), retain
\[
 n_3=\sum_{\substack{\ell\mid a\\\ell\equiv3\ (7)}}v_\ell,
 \qquad
 n_{24}=\sum_{\substack{\ell\mid a\\\ell\equiv2,4\ (7)}}v_\ell.
\]
Then \(E_a\cup M_a\ne\varnothing\) if and only if at least one
of the following three explicitly computed predicates holds:
\[
 \begin{gathered}
 \text{a prime divisor of }a\text{ is }5\text{ or }6\pmod7;\\
 n_3\ge3;\qquad n_3\ge1\text{ and }n_{24}\ge1.
 \end{gathered}
\]
The following choices give a witness in every passing case:
\[
 \begin{array}{c|c|c}
 \text{retained factors or divisor}&u&\text{channel}\\ \hline
 \ell\mid a,\ \ell\equiv5\ (7)&\ell&E_a\\
 \ell\mid a,\ \ell\equiv6\ (7)&a\ell&M_a\\
 d\mid a,\ d\text{ a product of three }3\ (7)\text{ prime factors}
        &ad&M_a\\
 \ell,t\mid a,\ \ell\equiv3\ (7),\ t\equiv2\ (7)
        &a\ell t&M_a\\
 \ell,t\mid a,\ \ell\equiv3\ (7),\ t\equiv4\ (7)
        &a\ell/t&M_a
 \end{array}
\]
The three factors defining \(d\) are counted with their original
prime multiplicities. No factor is used more often than it divides \(a\).

There is also an exact description of every witness, including those
not selected by this table. Set
\[
 \lambda(1)=0,\quad\lambda(2)=2,\quad\lambda(3)=1,\quad
 \lambda(4)=4,\quad\lambda(5)=5,\quad\lambda(6)=3.
\]
For each original prime \(\ell\mid a\), put
\(\lambda_\ell=\lambda(\operatorname{rem}_7\ell)\), and put
\(b=\sum_{\ell\mid a}\lambda_\ell v_\ell\) in \(\Z/6\Z\).
The full fibres, with the original valuation inverse, are
\[
 E_a=\{u(e):e\in\mathcal B_a,\ \sum_\ell\lambda_\ell e_\ell=5\ (6)\},
\]
\[
 M_a=\{u(e):e\in\mathcal B_a,\ \sum_\ell\lambda_\ell e_\ell=b+3\ (6)\}.
\]
In the finite group ring \(\Z[T]/(T^6-1)\), form
\[
 F_a(T)=\prod_{\ell\mid a}\sum_{e=0}^{2v_\ell}T^{\lambda_\ell e}
        =\sum_{r=0}^5 c_rT^r.
\]
Then \(|E_a|=c_5\), \(|M_a|=c_{\operatorname{rem}_6(b+3)}\).
Every coefficient counts its whole displayed exponent fibre; no
exponent representative replaces the finite integer budget.
\end{theorem}
\begin{proof}
The original unit calculation gives \(\gcd(a,7)=1\).
Explicitly, if seven divided \(a\), then \(p=4a-7\) would be
divisible by seven, contradicting primality and \(p\ge13\).
The powers \(3^0,\ldots,3^5\) modulo seven are
\(1,3,2,6,4,5\); the next power is one, and these six residues
are distinct. They therefore enumerate every nonzero residue
with the listed inverse \(\lambda\).
Since \(4^{-1}=2\pmod7\), the exterior target is
\(-4^{-1}=5=3^5\). The middle equation is \(u/a=-1=3^3\)
in the unit group. Taking these exact finite-group coordinates
proves the two displayed fibres. Expanding the finite product
\(F_a\) chooses one original exponent at each prime; collecting
terms with equal residues modulo six proves each count and its
entire underlying fibre.

Each table entry is a positive divisor of \(a^2\). For \(u=\ell\),
the exponent is one at \(\ell\) and zero elsewhere. For \(u=a\ell\)
the exponent at \(\ell\) increases from \(v_\ell\) to
\(v_\ell+1\le2v_\ell\). If \(n_3\ge3\), select exactly three
factors from the multiset having \(v_\ell\) copies of each
\(3\pmod7\) prime. Their product \(d\) divides \(a\), so
each exponent in \(ad\) is between \(v_\ell\) and \(2v_\ell\).
For the last two rows \(\ell\ne t\), since their residues differ.
For \(a\ell t\), both corresponding exponents increase by one
within their budgets. For \(a\ell/t\), the \(\ell\)-exponent
increases by one and the \(t\)-exponent decreases by one; the
latter is nonnegative since \(t\mid a\). All other exponents
remain exactly those of \(a\).

The first entry has residue five, hence belongs to \(E_a\).
The ratios \(u/a\) of the other entries have residues, in order,
\[
 6,\qquad 3^3=6,\qquad3\cdot2=6,\qquad3/4=3\cdot2=6
 \quad\pmod7.
\]
Thus they belong to \(M_a\). This proves sufficiency with full
integer and exponent checks.

For necessity suppose all three predicates fail. There is no prime
factor with residue five or six. If \(n_3=0\), every prime factor
has residue in \(H=\{1,2,4\}\). Direct multiplication modulo
seven shows \(H\) is the subgroup generated by two, of order
three. Every \(u\mid a^2\) and \(a\) itself then lie in \(H\).
The exterior target five is outside \(H\), and the middle target
\(-a\) is outside \(H\): otherwise its quotient by \(a\)
would put \(-1=6\) in \(H\). Both channels are empty.

The only remaining failed-predicate case has \(n_3=1\) or two
and \(n_{24}=0\). Every prime factor now has residue one or
three. Write \(e=\sum_{\ell\equiv3\ (7)}e_\ell\) for the
sum of the actual divisor exponents. Its range is
\(0\le e\le2n_3\le4\), and \(u\equiv3^e\pmod7\).
No such exponent is five modulo six, so the exterior target
is absent. For the middle condition,
\(u/a\equiv3^{e-n_3}\), with \(-2\le e-n_3\le2\).
None of these integers is three modulo six, so the middle
target is absent. This proves necessity without enlarging
the exponent intervals to their generated group.
\end{proof}

\begin{corollary}[The exact two-shell rejection set and integer reconstruction]
\label{cor:two-shell-sieve-return}
For every original prime \(p=12h+1\), the union of the two complete
shell loci is empty exactly when both of the following hold:
\[
 \begin{gathered}
 \text{every prime divisor of }3h+1\text{ is }1\pmod3;\\
 \text{no prime divisor of }3h+2\text{ is }5\text{ or }6\pmod7,\qquad
 n_3=0\ \text{or}\ (1\le n_3\le2\text{ and }n_{24}=0).
 \end{gathered}
\]
The counts and complete exponent fibres in the preceding theorems
give a sieve for these actual two shells for every such prime,
without a numerical cutoff on \(p\). A surviving prime is one
for which these two shells both fail; no assertion about the other
original shells, or the number of surviving primes, is implied.

For every retained passing pair \((\varepsilon,u)\), with
\(\varepsilon=0\) for the exterior channel and one for the middle,
set
\[
 g=\gcd(a,u),\quad A=a/g,\quad B=u/g,\quad D=g^2/u,\quad
 C=(A+p^{1-\varepsilon}B)/R_a.
\]
These are positive integers and give the exact original ordered witness
\[
 (a,y,z)=(ABD,p^\varepsilon ACD,pBCD),\qquad
 R_a=4a-p,\quad S_a=pa.
\]
The map from the retained tagged exponent vector to
\((p,a,R_a,S_a,\varepsilon,u,A,B,C,D,y,z)\) has singleton fibres
and inverse the original prime valuations of \(u\), retaining
\(p,a,\varepsilon\). All constructed witnesses have
\(a\le(p-1)/2\); membership in a smaller unary layer is imposed
on the computed \(A\), not assumed for every table branch.
\end{corollary}
\begin{proof}
Taking the conjunction of the two proved empty-fibre conditions
gives exactly the displayed rejection set, in both directions.
Each full locus remains available through its exponent fibre,
so the rejection calculation has not discarded any unselected
witness or any original prime factor.

To prove the stated reconstruction directly, fix a prime \(\ell\)
and write \(v=v_\ell(a)\), \(e=v_\ell(u)\), with \(0\le e\le2v\).
Then \(v_\ell(g)=\min(v,e)\), and the exponent of \(g^2/u\)
is \(2\min(v,e)-e\), which is \(e\ge0\) if \(e\le v\),
and \(2v-e\ge0\) if \(e\ge v\). Thus \(D\) is a positive
integer. The exact definitions give \(\gcd(A,B)=1\),
\(a=ABD\), \(u=B^2D\); their inverses are the displayed
formulas for \(g,A,B,D\).

All of \(A,B,D\) are units modulo \(R_a\), because their
prime factors divide \(a\), which is coprime to \(R_a\).
For the exterior channel multiply \(4u+1\equiv0\) by \(A\):
\[
 A(4B^2D+1)=4aB+A\equiv pB+A\pmod{R_a}.
\]
Thus \(R_a\mid A+pB\). For the middle channel,
\(u+a=BD(A+B)\), and the unit \(BD\) gives
\(R_a\mid A+B\). These prove that \(C\) is a positive integer
in each case, with no missing divisibility condition.
Now \((4a-p)C=A+p^{1-\varepsilon}B\), hence
\[
 4ABCD=A+p^{1-\varepsilon}B+pC.
\]
The three reciprocals, over the original common denominator
\(pABCD\), have numerators \(pC,p^{1-\varepsilon}B,A\).
Their sum is therefore \(4/p\), with precisely the asserted
ordering and unchanged \(a,R_a,S_a\). Conversely the graph
retains \(u\), so its original exponent vector and tag are
recovered exactly and uniquely. Finally
\(3h+1\le6h\) and \(3h+2\le6h\) for \(h\ge1\), proving
the retained first-half bound. The integer \(A=a/g\) remains
an explicit coordinate for any further layer condition.
\end{proof}

\paragraph{An occupied residual-seven shell outside the first two layers.}
The prime \(p=373=12\cdot31+1\) has
\(a_7=95=5\cdot19\). Primality is proved by the nonzero remainders
\(1,1,3,2,10,9,16,12\) on division by, respectively,
\(2,3,5,7,11,13,17,19\), all the primes at most
\(\sqrt{373}<20\). Both factors of \(a_7\) are five modulo seven.
Every original divisor is uniquely \(u=5^i19^j\),
\(0\le i,j\le2\), and its residue is \(5^{i+j}\).
For \(i+j=0,1,2,3,4\), these residues are
\(1,5,4,6,2\). The exterior target five occurs exactly at
\((i,j)=(1,0),(0,1)\), so \(E_{95}=\{5,19\}\).
The middle target is \(-95\equiv3\pmod7\), which does not occur;
thus \(M_{95}=\varnothing\).
Their complete primitive coordinates are, respectively,
\[
 (A,B,C,D,\varepsilon)=(19,1,56,5,0),\qquad
 (5,1,54,19,0),
\]
with ordered integer witnesses
\[
 (95,5320,104440),\qquad(95,5130,382698).
\]
Indeed \(7\cdot56=19+373\), \(7\cdot54=5+373\),
and each has \(ABD=95\); the preceding reciprocal calculation
therefore proves both identities with the original prime.
The full fibres show that every hit in this shell has \(A=19\)
or five. In particular no \(A\le2\) code in this shell is a hit.
This determines precisely the failure of the same-shell implication
from occupancy to the first two layers; all other original shells
remain in the full primewise system.
